% Partial authored draft only. Not included in review or counted as rewritten.
\chapter{表示距离：统计区分、路径代价与几何模型}
\label{chp:riemannian_base}

前一章讨论了关系基底、局部属性空间和统计区分度。本章进一步追问：系统说两个对象接近，究竟是在比较属性、统计分布、访问路径，还是实现成本？这些比较可以共同参与决策，却不能共享一个未经定义的距离。我们将先规定比较对象与有效域，再分别建立对称距离、统计尺度及路径模型，说明它们何时能够联合，何时必须保持分离。由此形成的几何不是预先给定的全部现实，而是带来源、版本与任务条件的计算模型。

我们还需要区分模型的稳定性和环境事实的约束。目标变化不能自行改变观测证据，事实记录也不能因被称为客观而免于修订。可达路径不等于有效推理，最短路径不等于最正确结论，保范数变换也不自动保证逻辑规则。我们将在固定结构条件下推导路径与局部度量的关系，再把接口更新、群变换和实现预算放回模型。HSF-HD提供这些一般状态动力学条件，AgentOS则在记忆、控制和边界中实例化它们，二者不因共用几何术语而成为同一个理论层次。

\section{比较对象与距离公理：对称性不提供因果方向}
\label{sec:section_203}

我们首先声明距离比较的对象。两个对象记录可以因为颜色接近而具有较小属性距离，也可以因为任务角色不同而具有较大的转换成本。两段程序可以输出相同结果却需要不同运行时间，两条数据库记录可以指向同一对象却具有不同来源。若不明确比较空间，接近一词就会把数值相似、任务等价和执行可达混在一起。我们把对象身份、属性编码、关系位置、证据分布及实现成本分别建模，再通过声明的读出接口使用这些量，而不设想一种距离独立决定对象的全部语义。

在集合 $X$ 上，一个距离函数 $d:X\times X\to[0,\infty)$ 需要非负性、零距离当且仅当对象相同、对称性及三角不等式。若零距离可以对应两个不同对象，我们得到的是伪度量，可以在零距离等价类上再研究商空间。若访问代价具有方向，去往目标和返回起点可能不同，它就不能未经处理成为对称距离。若允许不可达，我们可以使用扩展值或显式不可达标记，但算法必须说明如何处理这些值。任务把两个对象视为等价时，零距离也应说明是在任务商空间中相同，而不是现实身份已经合并。

固定有限维实属性空间中，设坐标按参考尺度归一化，$G$ 为对称正定矩阵。我们定义
\[
 d_G(x,y)=\sqrt{(x-y)^{\mathsf T}G(x-y)}.
\]
由正定性可取可逆矩阵 $C$ 使 $G=C^{\mathsf T}C$，于是距离等于 $\|C(x-y)\|_2$。非负性和对称性直接成立，可逆性使零距离只在两点相同时成立，欧氏范数的三角不等式则给出三角关系。因此我们得到一个真正的距离，而不只是把某个相似评分改名为距离。若矩阵仅半正定，零方向会压缩不同输入，得到的只是伪度量；若矩阵不对称，二次型只看到对称部分，不能借此表达所有方向依赖。

这个构造依赖固定比较规则。若我们对每一对输入临时选择不同的矩阵，对称性和三角不等式就需要重新检查；若矩阵随位置变化，则两点距离通常需要由路径长度的下确界定义，而不是只在一个端点套用固定二次型。在连续几何中，局部正定度量描述小变化的长度，整体距离还依赖空间连通性及允许路径。离散图中的跳数和属性空间中的二次距离也不是天然同一量：前者受边集合限制，后者可能允许连接在图上并不合法的对象。我们只在明确映射及误差条件时联合它们。

对称性尤其不能替代因果或逻辑方向。若某条件是另一事件的必要条件，逻辑上得到的是事件成立蕴含条件成立，不能据此倒转为条件成立就保证事件发生。两个变量统计相关，也不说明干预其中一个必然改变另一个。我们可以在无向层保存相似性，在有向层保存依赖或因果假设，在规则层保存蕴含条件；所有层共享身份与来源，却分别进行验证。一项对称距离小，只说明规定比较下接近，不说明两个命题相互推出，也不说明对象能够互换执行。

火与热的例子可用于说明建模粒度。我们需要先规定变量是某种测量、某类事件标签还是语词出现，并说明采样环境与条件。词语共现所得统计关系和环境测量所得关系不能直接视为同一对象，相关尺度也可能随语境和采样方案改变。把一项关联写入当前模型不会使它成为全宇宙固定距离。模型可以保持对证据的客观约束，同时承认观测不完整及参数可修订；客观性应体现为来源可核查和规则一致，而不是把矩阵数值宣布为不受任何更新影响。

猫和狗在某种哺乳动物属性编码中可能较近，在具体识别任务中却必须保持不同标签。对数函数和香草冰淇淋在一种字段空间中可能缺少直接可比属性，在教学或菜单定价任务中又可以具有有用关系。属性正交不表示图结构一定不连通，数值距离较大也不表示完全没有语义关联。我们让任务提供比较准则、适用范围和读出目标，保留不同准则之间的不一致作为信息，而不强制所有对象进入一个统一远近序列。

结构编辑代价是另一个需要单独定义的对象。若编辑操作及其逆操作具有相同成本，并且操作能连接全部比较对象，最小累计编辑成本可以产生对称路径距离；若操作不允许逆转，或合并与拆分代价不同，通常得到方向代价。零成本编辑还可能使不同记录处于同一伪度量类。我们在接口中列出操作、成本、合法条件和对象等价准则，再分析路径性质。这样，距离公理成为可检查的模型责任，而不是从一种本体论名称自动取得的保证。

\section{统计距离：互信息、变异信息与局部参数几何}
\label{sec:info_derivation_of_metric}

统计量比较的是随机变量在指定样本分布中的关系，而不是两个概念脱离数据后的永恒距离。我们设 $X$ 和 $Y$ 为有限集合上的离散随机变量，联合概率 $p(x,y)$ 非负且总和为一，边缘概率由求和得到，所有对数采用同一底数，零概率项按连续延拓记为零。熵、条件熵和互信息分别定义为
\[
H(X)=-\sum_xp(x)\log p(x),\qquad H(X\mid Y)=-\sum_{x,y}p(x,y)\log p(x\mid y),
\]
\[
I(X;Y)=\sum_{x,y}p(x,y)\log\frac{p(x,y)}{p(x)p(y)}.
\]
这些量按所选对数底数以比特或奈特计量，与状态活动、目标代价和实际能耗不同，不能把信息量的增减直接换算为运行功率或理解强度。

由 $p(x\mid y)=p(x,y)/p(y)$ 展开求和，可得 $H(X\mid Y)=H(X,Y)-H(Y)$，交换变量也有对应关系。把互信息的对数比拆成三项并按边缘求和，得到 $I(X;Y)=H(X)+H(Y)-H(X,Y)$。因此变异信息满足
\[
\operatorname{VI}(X,Y)=H(X\mid Y)+H(Y\mid X)=H(X,Y)-I(X;Y).
\]
这个恒等式来自同一联合分布，不能把不同数据集、不同采样窗口或不同版本的熵项任意拼接。VI对随机变量生成的有限划分构成距离，但若两个变量以概率一互相确定而标签不同，其值仍为零；严格比较的对象是随机划分的等价类，不是现实身份。

原表达 $1-I(X;Y)/H(X,Y)$ 在联合熵为正时等于 $\operatorname{VI}(X,Y)/H(X,Y)$，可用作零到一之间的无量纲差异尺度。联合熵为零时，两变量均为确定常量，分母为零，必须单独定义退化情形。每一对变量若用各自联合熵归一化，还可能改变三元组之间的比例关系；在调用三角不等式或最短路径算法前，需要对具体形式证明。没有证明时，我们称其为相对差异指标，而不称为全局测地距离。

互信息是对称依赖量，不提供因果方向。火与温度在某个数据集中高度相关，只说明该采样条件下共同不确定性减少；更换环境、测量窗口或控制变量后，数值可以变化。估计还受样本量、分箱、模型假设和重复来源影响。我们要求保存数据版本、估计器与不确定范围，把统计相关、机制假设和干预结果分开。高互信息可以生成候选关系，却不能把一条边自动升级为因果或逻辑必然边。

局部参数几何处理的是另一类对象。对分布族 $p_\theta(z)$，若支持集固定、概率为正且足够光滑，Fisher矩阵由评分外积期望定义。对小扰动 $\delta$，在交换求导与求和合法时有
\[
D_{\mathrm{KL}}(p_\theta\Vert p_{\theta+\delta})=\tfrac12\delta^{\mathsf T}I_F(\theta)\delta+o(\|\delta\|^2).
\]
一阶项因评分均值为零而消失，二阶项由归一化恒等式得到。该式只是同一模型族参数邻域中的二阶近似，不说明远距离KL对称，也不说明沿Fisher度量积分所得距离等于任意变量的VI。参数不可辨识时矩阵退化，只能得到伪长度，需要在等价类上分析或采用有说明的正则化。

我们因此保留三种比较：字段依赖使用互信息，随机划分差异使用VI，可微模型的局部参数可辨识性使用Fisher几何。它们有公式联系，却作用于不同对象和尺度。若统计差异成为关系图边权，还应说明权重如何随新数据更新、不可达怎样处理以及路径组合是否有意义。在AgentOS中这些尺度可辅助 $h(t)$ 的候选检索、$E$ 的事件聚类和 $\Theta$ 的结构抽象，但每次写入仍携带来源和任务条件；自我相关优先级能改变计算预算，却不能改写观测概率。验证时还应构造独立、完全相关、标签置换和样本稀疏等基准情形，确认估计器在已知边界下给出预期结果；只有通过这些局部检验，统计尺度才可进入后续控制，而不能因公式形式优美就被赋予普遍本体地位。

\section{世界结构快照：证据约束、相似网络与版本更新}
\label{sec:section_220}

世界结构是系统在某一版本上对对象、关系和规则的可检查承诺。我们称它为背景，是因为一次局部推理或控制步骤先固定其版本，使读出具有稳定参照，而不是认为它绝对静止。结构中可包含物理规律模型、数学公理、社会规则和经验常识，但验证方式不同：公理属于形式系统约定，物理模型依赖测量范围，社会规则依赖制度与时间，常识可能只在常见情境成立。它们处于同一图中并不获得相同真值地位。

我们以带类型多层图表示快照 $G^{(v)}=(V,E,\tau,\rho,\mathcal P)$。$v$ 是版本号，$\tau$ 给出节点和边的类型，$\rho$ 保存参数或规则，$\mathcal P$ 保存来源、时间和适用域。身份相同的节点可跨版本持续，属性和边通过事件日志增删。推理事务读取固定版本；若其间结构更新，提交时检查依赖是否变化，受影响结果重新计算或标记过期。这样的刚性来自版本一致性和证据约束，不是数值结构永远不可修改。

目标和情绪可以改变检索优先级、行动效用与资源分配，却不能单独改变已记录的外部测量。主体不喜欢重力不能让物体下落观测从日志消失；“跳下必然死亡”仍是带条件的风险模型，受高度、防护、环境及主体状态影响。我们将观测、预测和规范承诺分开：观测描述发生了什么，预测模型给出条件结果，规范约束决定哪些动作不可接受。偏好可以选择不执行某条路径，却不能充当证据否定环境反馈。

电阻网络可以作为有限图上的专门模型。设无向连通图边电导 $c_{ij}=c_{ji}>0$，拉普拉斯为 $L$。在节点 $a$ 注入单位流、节点 $b$ 抽取单位流并固定参考电位后，有 $Lu=e_a-e_b$，有效电阻为
\[
R_{ab}=(e_a-e_b)^{\mathsf T}L^+(e_a-e_b),
\]
其中 $L^+$ 是伪逆。该量对称、非负并综合并行路径，不等于单边阻值或最短路径。用于相似网络前要说明电导如何从证据得到；有向图、负权图和时变图不能直接套用标准解释。我们用它分析冗余可达性，而不把全部语义度量等同为实际电阻。

信息沿图传播依赖更新规则。无输入扩散可以偏向高电导路径，但联想生成不自动保证真值。大量重复网页形成的强边可能只是同源复制，低频边也可能是关键安全规则。结构写入时应进行来源去重、证据分级和反例检查；读出时区分候选生成与结论验证。相似网络负责找到可能相关的对象，规则和证据层负责检查能否组合。若压成一个权重，就难以识别高相似但逻辑不相容的错误。

结构更新需要变更集、前置条件和回滚范围。添加边前确认端点身份及关系类型，删除边时检查依赖结果，合并节点时保留别名和冲突字段，拆分节点时重新分配来源。内部记录可回滚，已经执行的现实动作只能通过补偿处理。结构变化会使旧距离、缓存和路径失效，因此都要绑定版本。跨版本迁移还应报告哪些指标可比较，不能把参数更新造成的距离变化误解为对象本身突然改变。

事实层也可能受观测误差和分布变化影响。我们允许记录置信范围与未决状态，并用新测量修订；可修订性不等于任意性，因为修订必须携带证据或规则变更。冲突来源可以作为多个候选解释并存，而非强行平均成一个精确结论。AgentOS可以把当前快照载入 $h(t)$，把变更事件写入 $E$，把反复验证的规则沉淀到 $\Theta$，完整版本放在外部记忆；一般HSF-HD只要求声明版本、来源和更新机制，不把世界图规定为不可变宇宙背景。

\section{路径优化与测地模型：驻值、最短与有效推理}
\label{sec:section_6308}

当状态空间可用连续坐标近似时，局部度量能够定义路径长度。设 $\mathcal M$ 为光滑流形，坐标按参考尺度归一化，$g(x)$ 为连续可微、对称正定的无量纲矩阵场。分段光滑路径 $x:[0,1]\to\mathcal M$ 的长度和能量为
\[
\mathcal L[x]=\int_0^1\sqrt{\dot x^{\mathsf T}g(x)\dot x}\,ds,\qquad
\mathcal E[x]=\tfrac12\int_0^1\dot x^{\mathsf T}g(x)\dot x\,ds.
\]
$s$ 是无量纲路径参数，不是物理时间或记忆内化时间。固定端点下，距离定义为允许路径长度的下确界；是否存在达到下确界的路径，还需要完备性或紧致性等条件。

为得到驻值方程，我们对能量作变分。取端点为零的扰动 $\eta(s)$，令 $x_\varepsilon=x+\varepsilon\eta$。对被积函数求导，对速度项分部积分，端点项因扰动为零消失。任意扰动下的一阶变化为零，得到
\[
\ddot x^\mu+\Gamma^\mu_{\nu\lambda}(x)\dot x^\nu\dot x^\lambda=0.
\]
Christoffel符号由度量和一阶偏导组合。该方程只是固定端点能量驻值的必要条件；局部最短还需检查二阶变化，全局最短更不能只凭微分方程断言。恒速参数化时长度驻值和能量驻值可对应，任意参数化下不能混用。

把该路径称为低偏置路径，只表示固定度量和允许集合中没有额外目标项。它不是无需计算代价的惯性滑行，也不保证每一步逻辑有效。若坐标只表达语义相似，最短路径可能穿过缺少规则支持的区域；若硬约束未写入允许集合，连续曲线还会跨过不可执行边界。几何模型用于候选路径生成，证明检查器、模拟器和环境反馈用于判定转换是否合法。

加入目标、风险和资源后，路径目标应明确改变。控制路径若同时产生任务代价 $c(x,u,t)$、风险约束与外部输入，最优控制问题不再等同于原度量测地线。目标随时间变化时需要在线重规划；结构或度量变化时，旧连接系数也应更新。不能令所谓规范场为零就声称消除了全部偏差，因为训练数据、坐标和度量估计本身包含建模选择。客观推理依靠来源、规则和反例检验，而不是只依赖一条平滑曲线。

离散实现要声明步长和数值方法。显式欧拉积分不会精确保持连续模型的速度范数，大步长还可能离开坐标有效域。可以采用更稳定的积分器、投影回合法集合，或者直接在图上求满足类型约束的路径。任何实现都要检查离散误差、停止条件和不可达状态。连续模型的性质不由任意程序自动继承，路径搜索结束也不意味着现实任务已经成功。

例如，从患者发热到某种药物的语义路径可能很短，因为文本中频繁关联；有效决策还需要年龄、过敏史、病因、剂量和授权。几何检索可迅速找到候选，规则与环境接口必须阻止未经验证的执行。数学证明中，每一步也需满足推理规则，不能用向量连续性补足缺失前提。最短路径只优化已经定义的长度，不负责定义正确性。

路径模型仍可用于比较表示搜索难度、定位高代价接口和分析局部扰动。多条等长路径存在时可保留多个候选；度量只在局部可靠时可分段重估并记录版本。在AgentOS中，$h(t)$ 承载候选和验证状态，$E$ 记录尝试与反馈，$\Theta$ 提供转换规则，TCE调节搜索尺度，CCM控制分支数，Boundary约束外部动作。目标下降只说明内部候选改善，必须经过现实反馈才判断成功。若反馈与内部预测冲突，系统应保留失败轨迹、重新估计模型并撤销尚未执行的后继计划，而不是调整距离定义去掩盖错误；这种可纠错性比路径是否光滑更能说明模型的工程价值。

\section{抽象计算与载体实现：内在规则、嵌入接口和资源预算}
\label{sec:intrinsic_vs_extrinsic_geometry}

我们区分抽象状态关系与载体实现，使同一计算模型能在不同硬件、软件和组织中讨论。内在描述规定对象、允许变换、距离和约束；外在实现规定这些对象怎样编码到内存、网络、传感器和执行器。抽象模型可以不依赖某块芯片的坐标，但现实系统需要时间、存储、通信和实际能耗。两层不能混同，也不能断开：实现必须保持任务读出，抽象设计也应接受资源预算和故障模式反馈。

设抽象状态集合为 $X$，实现状态集合为 $Z$，编码和读出分别为 $e:X\to Z$ 与 $r:Z\to X$ 或任务等价类。精确实现要求有效域内 $r(e(x))=x$；近似实现给出任务误差界。抽象更新 $F$ 与实现更新 $\widehat F$ 的一致性，是 $r(\widehat F(e(x)))$ 与 $F(x)$ 在指定度量下接近。这个关系需逐类操作验证，不能由可视化嵌入相似证明。实现状态还含缓存、调度和校验等辅助变量，不一定与抽象模型一一对应。

几何嵌入定理必须按条件使用。一般黎曼流形在足够高维欧氏空间可有等距嵌入，并不表示高维流形能等距塞入二维或三维空间，也不推出低维实现必然通过剧烈折叠保持全部度量。数字计算机可用地址、时间复用和协议模拟抽象邻接，无需物理相邻精确复制每条边。皮层沟回增加表面积并受发育与连接约束影响，不能从抽象定理直接推导其唯一功能机制。

物理布局仍影响性能。二维芯片、三维堆叠、光互连、存内计算和分布式节点在带宽、延迟、热管理、良率及故障隔离上有不同取舍。三维堆叠能缩短某些连接，也可能增加散热难度；高带宽能扩大并行交换，却不保证消息身份和来源正确。实际延迟按秒、存储按字节、吞吐按单位时间消息、能耗按焦耳计量，不能用抽象向量模长或活动强度代替。

逻辑约束也不能仅靠物理连线刻蚀。硬件可实现访问控制、类型检查和冗余校验，使某些非法操作更难发生；规则更新、异常输入和软件错误仍需系统处理。大语言模型的幻觉可能来自证据缺失、目标错配、解码策略、训练分布及工具失败，不能都归因于高维欧氏表示缺少一种壁垒。显式结构和验证器能降低某些错误，却也可能写入错误规则或过时事实。可靠性来自多层约束、反馈和审计。

局部结构网络与全局分布表示可以采用不同存储路径。图记录保存身份、关系和版本，分布表示支持相似检索和模式泛化；绑定接口将向量候选读出到明确对象，并把验证结果反馈给索引。缓存要绑定版本以防陈旧。外部工具、数据库和其他Agent进入闭环时，通信与权限边界也进入资源和安全模型。系统的有效身体由可感知、可行动和可维护的接口共同界定，而不只是参数规模或芯片形状。

我们还区分可逆内部步骤与现实不可逆后果。程序可复制状态、回滚事务或重算路径，外部动作却可能消耗材料、伤害对象或改变承诺。若把所有操作当成坐标变换，会低估执行成本和责任。Boundary应在动作前检查授权与可逆性，情景记忆记录结果，必要时触发补偿。AgentOS通过 $h(t)$、$E$、$\Theta$ 和外部记忆协调抽象—实现契约，CCM管理复杂度，Offloading重分配存储；这是一套工程方案，HSF-HD仍对其他实现开放。

\section{变换与更新：保范数条件、尺度变化和路由验证}
\label{sec:group_theory_morpho_semantics}

不同变换保持不同性质。实向量 $x\in\mathbb R^d$ 使用欧氏范数，线性变换 $Q$ 对所有向量保范数，当且仅当 $Q^{\mathsf T}Q=I$；复空间的条件为 $U^\dagger U=I$。充分性由代入得到，必要性由所有向量的二次型相等推出矩阵为零。正交或酉变换保持内积和维数，却不自动保持对象身份、逻辑真值或任务约束，因为这些还取决于读出和类型接口。

一般可逆线性变换可旋转、剪切和缩放，非可逆映射还能压缩维度。我们不能规定形只能由正交群变换、质必须由缩放群变换：结构关系会通过增删边发生不可逆变化，属性内容也可在坐标重标记中保范数。更稳妥的做法是为每项操作声明不变量：身份重编码保持任务读出，属性校准保持测量含义，结构编辑保持引用完整性，压缩给出误差界。群结构是分析工具，不是两类语义永久分隔的自然法则。

非零复向量可写成 $x=ru$，其中 $r=\|x\|>0$、$\|u\|=1$；零向量没有唯一方向。若再把单位方向写成统一相位与内部方向的乘积，会出现规范不唯一：相位可以在两部分间转移而原向量不变。实高维向量也不能普遍由单个角度表示。因此径向代表内容、相位代表结构不是唯一分解定理，除非另行规定基底、规范和读出。我们保留幅度—方向分解分析活动尺度，但不把两部分直接等同于形和质。

连续更新 $\dot x=A(t)x$ 的范数导数为
\[
\frac d{dt}\|x\|^2=x^\dagger(A^\dagger+A)x.
\]
所以 $A$ 为斜Hermitian矩阵时保持范数；若含输入 $b(t)$，还要增加 $2\operatorname{Re}(x^\dagger b)$，不能沿用齐次结论。离散更新 $x_{k+1}=M_kx_k$ 对所有状态保范数，需 $M_k^\dagger M_k=I$。遗忘、门控、归一化和外部注入都会改变条件，数值范数改变只说明所选尺度变化，不等于焦耳能耗或感受强度。

路由常用归一化向量的余弦相似度，因为它消除整体尺度影响。这能减少大模长直接支配点积，却不能保证关系合法。不同对象可能方向相似，恶意输入也可刻意对齐，零向量需特殊处理。若关系具有方向、角色或时间条件，单一余弦值无法表达。我们让相似度负责候选召回，让类型掩码、身份绑定、来源检查和规则验证负责合法性；全部通过后，外部动作仍需授权和反馈。

注意力机制能表达统计关系，也可能错误聚合。错误不能只归因于幅度越权，因为归一化后仍可能缺证据，显式图也可能有错误边。我们分别测试尺度扰动、对象置换、无关高频项、来源冲突和逻辑反例，观察路由是否保持预期。归一化若提高稳健性，只能报告对应测试结果，不能宣称根除全部幻觉。

工程上可采用两阶段接口：先在全局分布表示中生成候选，再把候选绑定到局部结构记录并检查规则。每个候选保存评分、身份、来源与版本；没有合法对象时返回未决。CCM限制候选量，TCE调节探索，SPC压缩运行状态，Boundary控制执行。这个AgentOS方案不是一般智能唯一分解。本章结论是每个变换必须说明作用对象、保持量、允许损失及验证方法；下一章研究价值和控制如何改变选择方向时，也必须保持事实距离与目标偏好的区分。部署评估还要分别报告候选召回、绑定准确率、规则拒绝率、错误执行率和资源开销，避免用单一平均分掩盖安全边界上的失败。只有这些指标共同进入反馈闭环，路由机制的改进才会转化为可验证的系统可靠性。
